\subsection{Restitution fonctionnelle du moment orienté}
\label{iii:sec:restitucion-funcional-catalan}

La relation différentielle~\eqref{iii:eq:catalan-vacuum} admet un
inverse qui conserve la famille complète de moments. Son domaine est
la fonction constitutive produite par les secteurs du cycle, avec
l’orientation de~\eqref{iii:eq:uv}. Les incidences $90=3\cdot30$ et
$120=4\cdot30$ fixent cette fonction avant de l’évaluer dans les canaux
$s_\pm=q_\pm^{30}$. La condition à l’origine correspond à la
disparition du moment lorsque son argument contractant s’éteint.
Cette condition détermine également les constantes d’intégration.

\begin{theorem}[Restitution intégrale et unicité]
\label{iii:thm:restitucion-funcional-catalan}
Soit $f(s)=(1-s^3)/(1-s^4)$, pour $0<s<1$, la réponse
de~\eqref{iii:eq:trace-det}, et soit $\mathcal D=s\,d/ds$.
Il existe une unique fonction $F\in C^2((0,1))$ qui satisfait
\begin{equation}
 \mathcal D^2F(s)=\frac{s(1+s)}{1+s+s^2}f(s),
 \qquad \lim_{s\downarrow0}F(s)=0.
 \label{iii:eq:restitucion-funcional-problema}
\end{equation}
Elle est donnée par
\begin{align}
 F(s)&=\int_0^s\log\frac{s}{t}\,
          \frac{1+t}{1+t+t^2}f(t)\,dt
       =\int_0^s\frac{\log(s/t)}{1+t^2}\,dt
       =\mathcal C_2(s),
 \label{iii:eq:restitucion-funcional-integral}\\
 f(s)&=\frac{1+s+s^2}{s(1+s)}\mathcal D^2F(s).
 \label{iii:eq:restitucion-funcional-inversa}
\end{align}
L’extension continue à $s=1$ retrouve $F(1)=G_{\rm Cat}$.
\end{theorem}
\begin{proof}
La factorisation $1+t+t^2+t^3=(1+t)(1+t^2)$ donne
\[
 \frac{1+t}{1+t+t^2}f(t)=\frac1{1+t^2}.
\]
L’intégrande de~\eqref{iii:eq:restitucion-funcional-integral} est
intégrable en zéro, parce que
\[
 0\leq F(s)\leq\int_0^s\log(s/t)\,dt=s.
\]
Par conséquent, $F(s)\to0$. La différentiation sur l’intervalle ouvert,
d’abord sur un intervalle tronqué, puis par convergence dominée,
produit
\[
 \mathcal DF(s)=\int_0^s\frac{dt}{1+t^2},
 \qquad \mathcal D^2F(s)=\frac{s}{1+s^2}
     =\frac{s(1+s)}{1+s+s^2}f(s).
\]
En particulier, $F$ est de classe $C^2$ et satisfait l’équation.
Si $F_1,F_2$ sont deux solutions, la différence $H=F_1-F_2$
satisfait $\mathcal D^2H=0$. Au moyen de $z=\log s$, on obtient
$d^2H(e^z)/dz^2=0$, donc $H(s)=a\log s+b$.
L’existence de la limite nulle lorsque $s\downarrow0$ impose d’abord
$a=0$, puis $b=0$. La condition sur $F$ détermine donc
la solution et également la limite $\mathcal DF(s)\to0$.

Pour $s<1$, le développement géométrique de $(1+t^2)^{-1}$ et
l’intégrabilité absolue de ses termes pondérés permettent d’intégrer :
\[
 \int_0^s t^{2k}\log(s/t)\,dt
     =\frac{s^{2k+1}}{(2k+1)^2},\qquad
 F(s)=\sum_{k\geq0}\frac{(-1)^ks^{2k+1}}{(2k+1)^2}.
\]
C’est précisément la famille~\eqref{iii:eq:c2}. La convergence
absolue et uniforme de cette dernière série sur $[0,1]$ autorise le
passage à l’extrémité $s=1$. La borne
intégrable $\log(1/t)$ l’autorise également dans l’intégrale, en prolongeant l’intégrande par zéro pour
$t>s$. Enfin, isoler $f$ dans l’équation différentielle donne
\eqref{iii:eq:restitucion-funcional-inversa}, avec des dénominateurs
positifs sur $(0,1)$.
\end{proof}

La limite du moment et celle de sa dérivée seconde sont prises
dans leurs classes de convergence respectives. En particulier,
\[
 \lim_{s\uparrow1}\mathcal D^2\mathcal C_2(s)=\frac12
\]
est la limite radiale d’Abel de la série dérivée. L’évaluation
de $\mathcal C_2(1)$ utilise directement sa série absolument
convergente. Cette distinction conserve le même opérateur de profondeur
à l’intérieur et dans sa lecture de frontière.

\subsection{Compatibilité de la restitution avec les canaux et la vitesse}
\label{iii:sec:restitucion-canales-velocidad}

Le théorème restitue une fonction à partir d’une autre fonction préalablement
construite. L’inversion cubique~\eqref{iii:eq:cubic} agit, en revanche,
sur les évaluations $r_\pm$ : elle retrouve les arguments $s_\pm$
au sein de cette famille fixée. La composition des deux opérations
conserve l’orientation du quart de tour, les étiquettes $P_\pm$
et les incidences du transport. Ainsi, $G_{\rm Cat}=\mathcal C_2(1)$
est la lecture à l’extrémité d’un objet fonctionnel qui participe également
aux réponses intérieures $r_\pm=f(s_\pm)$.

Pour $\mathcal V_m=\mathcal V\otimes I_{9^m}$ et
$\tau_m=\operatorname{Tr}/(2\cdot9^m)$, la lecture de vitesse
maintient exactement la normalisation de
\eqref{iii:eq:velocidad-traza-normalizada} :
\begin{equation}
 \exp[-2\tau_m(\log\mathcal V_m)]
   =\frac1{r_+r_-}
   =(\det\mathcal V)^{-1}=\widehat c.
 \label{iii:eq:restitucion-velocidad-normalizada}
\end{equation}
En effet, chaque valeur propre $r_\pm$ se répète $9^m$ fois, de sorte
que $2\tau_m(\log\mathcal V_m)=\log r_++\log r_-$.
L’identité utilise le déterminant de la réalisation à deux feuillets
et la trace normalisée de son amplification. Les sections
dimensionnelles déjà fixées donnent alors
\[
 c_{\rm int}=c_{\rm vac}=\widehat c\,u_C,
 \qquad \ell_0=c_{\rm vac}t_0,
\]
avec le même état, la même coordinatisation et la même durée élémentaire
de~\eqref{iii:eq:longitud-elemental-constitutiva}.
La restitution fonctionnelle conserve cette section de vitesse lorsqu’on
la compose avec l’action et la longueur ; la multiplicité spectrale
est absorbée par la normalisation déclarée.
